\section{符号说明}\label{sec:notation}

下文只列出四问共用的核心符号；各问特有的情景索引、校准块编号、策略代号与算法中间量在相应小节首次出现时定义。

\begin{center}
\begin{tabularx}{0.96\linewidth}{@{}lYr@{}}
\toprule
符号 & 含义 & 单位 \\
\midrule
$d,t$ & 日期、时段索引 & -- \\
$p_{d,t}$ & 交付时段购电价 & 元/kWh \\
$L_{d,t},P_{d,t}$ & 负荷、光伏电量 & kWh \\
$\widehat L_{d,t},\widehat P_{d,t}$ & 负荷、光伏预测 & kWh \\
$E_{d,t}$ & 时段末储电量 & kWh \\
$c_{d,t},r_{d,t}$ & 充电量、放电量 & kWh \\
$w_{d,t}$ & 弃光量 & kWh \\
$\eta_c,\eta_d$ & 充、放电效率 & -- \\
\midrule
$q_{d,t}$ & 日初锁定的普通购电额度 & kWh \\
$x_{d,t}$ & 普通购电实际取用量 & kWh \\
$e_{d,t}$ & 紧急购电量 & kWh \\
$g^0_{d,t},g_{d,t}$ & 原计划、有效普通承诺 & kWh \\
$\phi_{d,t}(g^0,g)$ & 一次性调整结算费用 & 元 \\
$V_{d+1}(\cdot)$ & 日末储能机会成本 & 元 \\
\midrule
$C_d$ & 第 $d$ 日实际购电总成本 & 元 \\
\bottomrule
\end{tabularx}
\end{center}

\noindent 问题二至问题四的结果表只覆盖 2 月 1 日至 12 月 31 日，共 $334$ 日。1 月用于积累历史数据和预热 SOC，不写入题设工作簿；问题一则对应单日 $144$ 时段计划。文中成本均注明是输出区间还是 1--12 月连续运行口径。
